\chapter{The hardware bridge}
\label{ch:bridge}

\begin{safetybox}[Read this before the rest of the chapter]
Everything from here on can move a two-metre dual-arm robot. The operational
procedure --- network setup, sonar, enable, untuck, abort --- lives in the
repository's hardware runbook and command sheet, and assumes a supervised session
with the e-stop in hand. This chapter explains how the software works and why;
it is not a substitute for the runbook.
\end{safetybox}

\section{Shape of the problem}

Baxter runs a ROS~1 Noetic master. ROS~2 Jazzy cannot talk to it. The conventional
answer is \texttt{ros1\_bridge} on a host with both distributions installed, which
means maintaining a Noetic environment on Ubuntu 24.04 --- a container, in
practice, and a five-gigabyte one.

This workspace takes a different route. \texttt{scripts/py\_bridge.py} is a
pure-Python ROS~1 client: XML-RPC to the master, raw TCPROS to peers, no
\texttt{rospy} and no ROS~1 installation anywhere. It runs on the same laptop as
the ROS~2 side.

\begin{figure}[h]
\centering
\resizebox{\linewidth}{!}{% Hardware command path: ROS 1 master on the robot, through the pure-Python
% bridge, to the ROS 2 side. Drawn from scripts/py_bridge.py and
% src/baxter_hardware_bridge/.
\begin{tikzpicture}[
  font=\small,
  node distance=6mm,
  box/.style={draw, rounded corners=2pt, align=center, minimum height=8mm,
              inner xsep=5pt, fill=white},
  ros1/.style={box, fill=warnbg, draw=warnrule},
  ros2/.style={box, fill=notebg, draw=noterule},
  bridge/.style={box, fill=codebg, draw=coderule, very thick},
  flow/.style={-{Stealth[length=2mm]}, thick},
  lbl/.style={font=\scriptsize\ttfamily, inner sep=1.5pt, fill=white},
]

\node[ros1] (robot)  {Baxter\\\scriptsize ROS~1 Noetic master};
\node[ros1, below=8mm of robot] (jointstates) {\ttfamily/robot/joint\_states\\\ttfamily/robot/state};
\node[ros1, above=8mm of robot] (jointcmd) {\ttfamily/robot/limb/*/\\\ttfamily joint\_command};

\node[bridge, right=22mm of robot, minimum height=30mm, text width=24mm]
  (bridge) {\ttfamily py\_bridge.py\\[2pt]\scriptsize XML-RPC\\\scriptsize raw TCPROS\\[2pt]\scriptsize no rospy,\\\scriptsize no ROS~1 install};

\node[ros2, right=20mm of bridge, text width=26mm] (shim)
  {\scriptsize FollowJointTrajectory\\\small shim (per arm)\\[2pt]\scriptsize validate $\rightarrow$ safety gate\\\scriptsize $\rightarrow$ interpolate at 100\,Hz};

\node[ros2, above=10mm of shim, text width=26mm] (moveit) {MoveIt~2\\\scriptsize\ttfamily move\_group};
\node[ros2, below=10mm of shim, text width=26mm] (clients) {Example clients\\\scriptsize\ttfamily sim\_tiny\_trajectory};

\node[box, right=16mm of shim, draw=noterule, text width=20mm] (safety)
  {\ttfamily baxter\_\\safety\_check};

\draw[flow] (jointstates.east) -- ++(6mm,0) |- (bridge.west);
\draw[flow] (bridge.west) ++(0,6mm) -- ++(-6mm,0) |- (jointcmd.east);
\draw[flow] (bridge) -- node[lbl,above]{ROS~2} (shim);
\draw[flow] (shim) -- (bridge);
\draw[flow] (moveit) -- (shim);
\draw[flow] (clients) -- (shim);
\draw[flow] (shim) -- (safety);

\begin{scope}[on background layer]
  \node[draw=warnrule, dashed, rounded corners, inner sep=4mm,
        fit=(jointcmd)(robot)(jointstates), label={[font=\scriptsize\bfseries,
        text=warnrule]below:ROS 1 / robot}] {};
  \node[draw=noterule, dashed, rounded corners, inner sep=4mm,
        fit=(moveit)(shim)(clients)(safety), label={[font=\scriptsize\bfseries,
        text=noterule]below:ROS 2 Jazzy / laptop}] {};
\end{scope}

\end{tikzpicture}
}
\caption{The hardware command path. The bridge is a ROS~1 protocol client, not a
ROS~1 installation; the shim is where every safety decision is made.}
\label{fig:arch}
\end{figure}

\subsection{Why this is smaller than it sounds}

Subscribing to a ROS~1 topic requires negotiating a TCPROS connection, and the
handshake includes an MD5 sum of the message definition. The bridge answers
\texttt{"md5sum": "*"} --- a wildcard --- on the subscribe side. The master accepts
it, so \emph{receiving} a new message type needs no MD5 table and no message
definition, only code that can interpret the payload bytes. That is why adding a
topic is usually a dozen lines.

Publishing is not symmetric. \texttt{rosbag} records whatever the publisher
advertised, so a wildcard MD5 produced recordings that could not be deserialised
afterwards --- our own data, unreadable. The publisher side now answers with the
real MD5 and message definition for the types it publishes. Unlisted types still
fall back to the wildcard and still work on the wire; they just do not record
usefully.

\section{The action shim}

\texttt{FollowJointTrajectory} is what MoveIt and the example clients speak. The
robot does not: it accepts position commands on \texttt{joint\_command} and
expects them to keep arriving. The shim bridges that gap, one instance per arm,
and is where every safety decision lives.

\begin{figure}[h]
\centering
\resizebox{0.92\linewidth}{!}{% What the shim checks before a goal is allowed to move the arm.
% Transcribed from follow_joint_trajectory_shim.py; the 25 dry_run_test cases
% are the executable version of this diagram.
\begin{tikzpicture}[
  font=\small, node distance=5.5mm,
  stepbox/.style={draw, rounded corners=2pt, align=center, text width=52mm,
               minimum height=7mm, inner sep=3pt, fill=codebg, draw=coderule},
  dec/.style={draw, diamond, aspect=2.6, align=center, inner sep=1pt,
              text width=34mm, font=\footnotesize, fill=notebg, draw=noterule},
  bad/.style={draw, rounded corners=2pt, align=center, fill=warnbg,
              draw=warnrule, font=\footnotesize, text width=30mm, inner sep=3pt},
  ok/.style={draw, rounded corners=2pt, align=center, fill=white,
             draw=string, font=\footnotesize, text width=30mm, inner sep=3pt},
  flow/.style={-{Stealth[length=2mm]}, thick},
]

\node[stepbox] (goal) {Goal received};
\node[dec, below=of goal] (names)  {joint names\\match this arm?};
\node[dec, below=of names] (counts) {point sizes\\consistent?};
\node[dec, below=of counts] (limits) {positions inside\\URDF limits?};
\node[dec, below=of limits] (vel) {segment velocity\\within limits?};
\node[dec, below=of vel] (fresh) {joint states fresh\\($<2$\,s)?};
\node[dec, below=of fresh] (state) {exactly one\\\ttfamily/robot/state\\publisher?};
\node[dec, below=of state] (enabled) {robot enabled,\\not stopped, no error?};
\node[ok, below=of enabled] (accept) {ACCEPT\\\scriptsize interpolate at 100\,Hz};

\foreach \a/\b in {goal/names, names/counts, counts/limits, limits/vel,
                   vel/fresh, fresh/state, state/enabled, enabled/accept}
  \draw[flow] (\a) -- (\b);

\node[bad, right=14mm of names]   (r1) {REJECT\\\scriptsize bad joint names};
\node[bad, right=14mm of counts]  (r2) {REJECT\\\scriptsize malformed point};
\node[bad, right=14mm of limits]  (r3) {REJECT\\\scriptsize out of limits};
\node[bad, right=14mm of vel]     (r4) {REJECT\\\scriptsize over velocity limit};
\node[bad, right=14mm of fresh]   (r5) {REJECT\\\scriptsize stale state};
\node[bad, right=14mm of state]   (r6) {REJECT\\\scriptsize contested state};
\node[bad, right=14mm of enabled] (r7) {REJECT\\\scriptsize safety gate};

\foreach \a/\b in {names/r1, counts/r2, limits/r3, vel/r4, fresh/r5,
                   state/r6, enabled/r7}
  \draw[flow] (\a) -- node[font=\scriptsize, above]{no} (\b);

\end{tikzpicture}
}
\caption{What a goal has to pass before the arm moves. The 25 cases in
\texttt{dry\_run\_test} are the executable form of this diagram, and run in CI on
every push.}
\label{fig:gate}
\end{figure}

Once a goal is accepted, the shim interpolates it at 100\,Hz and publishes
position commands for as long as the goal is active. Because
\texttt{joint\_command} has a 0.2\,s timeout, holding a position means continuing
to command it --- so the shim keeps commanding the final point for
\texttt{hold\_duration\_sec} after the trajectory ends rather than falling silent
and letting the arm sag into gravity compensation.

\subsection{Tolerances}

Two checks run during motion, both taken from the values the legacy Rethink
trajectory action server actually shipped enabled
(\texttt{PositionJoint\-Trajectory\-Action\-Server}):

\begin{itemize}
  \item \texttt{path\_tolerance\_rad} (0.2): abort if a joint lags the command
    actually published by more than this. Note \emph{published}, not the raw
    interpolated setpoint --- comparing against a setpoint the shim never sent
    would blame the arm for the shim's own clamping.
  \item \texttt{stopped\_velocity\_tolerance} (0.25): after the final point has
    been held for \texttt{goal\_time\_sec}, abort if a joint is still moving
    faster than this.
\end{itemize}

Both \emph{hold} rather than stop, and they hold the \emph{measured} pose. An arm
lagging its setpoint is not necessarily an unsafe arm --- it may simply be blocked
--- and continuing to command a position it cannot reach would drive it harder
into whatever is blocking it.

\subsection{Distinguishing the failure modes}

\texttt{control\_msgs} has standard codes for tolerance failures
(\texttt{PATH\_TOLERANCE\_VIOLATED} $=-4$, \texttt{GOAL\_TOLERANCE\_VIOLATED}
$=-5$) and the shim uses them. It has no code for "cancelled" or "aborted by
safety gate", so the shim uses distinct negative values ($-100$, $-101$) rather
than leaving them at zero. A client that only reads \texttt{error\_code} must not
be able to mistake a safety abort for success.

\section{The safety gate}

\texttt{baxter\_safety\_check} reads \texttt{/robot/state} and reports whether
motion is permissible. One of its checks is less obvious than the rest: it
requires that \emph{exactly one} node publishes \texttt{/robot/state}.

That rule came from an incident. A mock robot left running from an earlier test
published a cheerful "safe" state alongside the real, disabled robot. A gate that
samples the latest message cannot tell which robot it is reading, and it went
unnoticed for about fifteen minutes. Counting publishers is the fix; the
hardware bringup launch also refuses to start if the mock is running.

\section{What still needs the legacy container}

Almost nothing. Bring-up, the interlock tests, supervised motion, tolerance work,
MoveIt and shutdown all run container-free. The exception is untucking.

Tucked arms sit outside the URDF limit table
(Section~\ref{ch:platform}), so the shim rejects that pose by design --- the same
guard rail that makes the workspace refuse to move while tucked. A container-free
untuck would mean hand-publishing joint commands \emph{outside the safety
envelope}, with collision avoidance suppressed, for a large whole-arm motion.

The honest answer is to keep the legacy tool for that one step. The alternative
--- relaxing the shim's limits behind an explicit opt-in so the untuck runs through
the validated path --- is the better long-term design, but it weakens a guard rail
that has already proved useful, and it must never be the default.
